\chapter{绪论}

\section{研究背景与意义}

电机是现代工业生产中应用最广泛的动力设备之一，广泛应用于电力、冶金、石化、制造等领域。据统计，全球工业用电量的70\%以上由电机消耗，电机的安全可靠运行直接关系到生产效率和设备安全\cite{nandi2005condition}。然而，电机在长期运行过程中，由于电气、机械、热力等多种因素的耦合作用，不可避免地会出现各种故障。常见的电机故障包括轴承故障、转子故障、定子故障和气隙偏心等，其中轴承故障约占电机故障总数的40\%$\sim$50\%\cite{zhang2020bearing}。

传统的电机故障诊断方法主要包括基于信号处理的方法和基于模型的方法。基于信号处理的方法通常需要先提取振动信号或电流信号的时域、频域或时频域特征，再利用分类器进行故障识别。这类方法严重依赖人工特征提取的经验和技巧，不同特征提取方法对诊断结果影响较大，且泛化能力有限\cite{lei2020applications}。基于模型的方法需要建立精确的电机数学模型，但在实际工业环境中，模型参数往往难以精确获取，限制了其工程应用。

近年来，深度学习技术在图像识别、自然语言处理等领域取得了突破性进展，为电机故障诊断提供了新的技术途径。深度学习能够自动从原始数据中学习多层次的特征表示，避免了人工特征提取的局限性，有望实现端到端的智能故障诊断\cite{zhao2019deep}。然而，将深度学习应用于电机故障诊断仍面临以下挑战：（1）工业现场故障样本稀缺，难以满足深度学习模型的大样本训练需求；（2）不同工况下信号特征差异大，模型的泛化能力有待提高；（3）深度学习模型的可解释性较差，难以满足工业应用的可信度要求。

因此，研究基于深度学习的电机故障诊断方法，对于提高电机故障诊断的准确性和智能化水平，保障工业生产的安全稳定运行，具有重要的理论意义和工程应用价值。

\section{国内外研究现状}

\subsection{基于传统方法的电机故障诊断研究}

传统的电机故障诊断方法主要分为基于信号处理的方法和基于模型的方法两大类。在信号处理方面，快速傅里叶变换（FFT）、小波变换、经验模态分解（EMD）等方法被广泛应用于电机故障特征提取。Nandi等\cite{nandi2005condition}对基于电流信号分析的电机故障诊断方法进行了系统综述，指出定子电流信号分析（MCSA）是最具工程应用前景的非侵入式诊断方法之一。

在模式识别方面，支持向量机（SVM）、随机森林（RF）、k近邻（KNN）等传统机器学习方法被广泛用于故障分类。Lei等\cite{lei2020applications}提出了基于集成经验模态分解和多尺度熵的轴承故障诊断方法，取得了较好的诊断效果。然而，这些方法都需要先进行人工特征提取，特征选择的好坏直接影响诊断性能。

\subsection{基于深度学习的电机故障诊断研究}

随着深度学习技术的发展，越来越多的研究者将其应用于电机故障诊断领域。Jia等\cite{jia2016deep}首次将深度神经网络（DNN）应用于旋转机械故障诊断，验证了深度学习在故障特征自动提取方面的优势。Zhang等\cite{zhang2018deep}提出了一种基于卷积神经网络（CNN）的轴承故障诊断方法，直接从原始振动信号中学习故障特征，实现了端到端的故障诊断。

在迁移学习方面，Li等\cite{li2020systematic}提出了基于迁移学习的跨工况轴承故障诊断方法，利用源域数据预训练模型，再在目标域上进行微调，有效解决了变工况下的故障诊断问题。Wen等\cite{wen2020new}提出了一种基于迁移学习和残差网络的故障诊断方法，在跨设备诊断任务中取得了较好的效果。

\subsection{现有研究的不足}

尽管基于深度学习的电机故障诊断方法取得了一定的进展，但仍存在以下不足：（1）大多数研究集中在公开数据集上的验证，缺乏在实际工业场景中的应用验证；（2）小样本条件下的诊断精度有待提高；（3）多故障耦合条件下的诊断能力不足；（4）模型的可解释性和可靠性有待增强。

\section{本文主要研究内容}

本文针对电机故障诊断中特征提取困难、小样本诊断精度低等问题，开展基于深度学习的电机故障诊断方法研究。主要研究内容包括：

\begin{enumerate}
	\item 电机故障振动信号分析。分析电机常见故障的振动信号特征，建立电机故障振动信号数据集，为后续研究提供数据基础。
	\item 基于CNN的端到端故障诊断方法。设计适用于一维振动信号的CNN网络结构，实现从原始信号到故障类别的端到端诊断。
	\item 基于迁移学习的小样本故障诊断方法。引入迁移学习策略，利用源域数据预训练模型，在目标域小样本条件下进行微调，提高诊断精度。
	\item 基于多尺度特征融合的故障诊断方法。设计多尺度特征融合网络，提取不同层次的故障特征，提高模型对多类型故障的识别能力。
\end{enumerate}

\section{本文组织结构}

本文共分为六章，各章内容安排如下：

第一章为绪论，介绍研究背景与意义，综述国内外研究现状，明确本文的研究内容和组织结构。

第二章为相关理论基础，介绍电机故障类型与机理、深度学习基础理论、卷积神经网络和迁移学习的基本原理。

第三章为基于CNN的端到端电机故障诊断方法，设计适用于振动信号的CNN网络结构，在公开数据集上验证方法的有效性。

第四章为基于迁移学习的小样本电机故障诊断方法，提出迁移学习策略，解决小样本条件下的诊断问题。

第五章为基于多尺度特征融合的电机故障诊断方法，设计多尺度特征融合网络，提高多类型故障的诊断精度。

第六章为总结与展望，总结本文的主要工作和创新点，分析存在的不足，并对未来研究方向进行展望。